\documentclass[11pt]{article}

\usepackage[margin=1in]{geometry}
\usepackage{array}
\usepackage{booktabs}
\usepackage{longtable}
\usepackage{tabularx}
\usepackage{xcolor}
\usepackage{hyperref}

\hypersetup{
  colorlinks=true,
  linkcolor=black,
  urlcolor=blue
}

\setlength{\parindent}{0pt}
\setlength{\parskip}{6pt}
\renewcommand{\arraystretch}{1.18}

\newcolumntype{Y}{>{\raggedright\arraybackslash}X}
\newcommand{\markcell}[2]{\textbf{#1/#2}}

\begin{document}

\begin{center}
  {\Large Geography Paper 2 HLSL 2022: Student Work Marking Report}\\[4pt]
  {\normalsize Marked against \texttt{Geography\_paper\_2\_\_HLSL\_markscheme-2022.pdf}}
\end{center}

\section*{Overall Result}

\begin{tabularx}{\textwidth}{@{}l c c Y@{}}
\toprule
Section & Marks awarded & Maximum & Comment \\
\midrule
Question 1: Changing population & 10 & 10 & Strong, concise answers. Physical and fertility explanations both developed clearly. \\
Question 2: Global climate & 10 & 10 & Accurate definitions and explanations. Migration answer was general but sufficient for the markscheme. \\
Question 3: Resource consumption and security & 9 & 10 & Good explanations; one mark lost for missing numerical support in the slum trend description. \\
Question 4: E-waste & 9 & 10 & Strong use of infographic evidence; one mark lost because the 6-mark evaluation needed an explicit overall judgement. \\
\midrule
\textbf{Total} & \textbf{38} & \textbf{40} & Excellent script; main improvements are quantifying graph trends and ending evaluation answers with a clear judgement. \\
\bottomrule
\end{tabularx}

\section*{Question-by-Question Marks}

\begin{longtable}{@{}p{0.12\textwidth} p{0.12\textwidth} p{0.31\textwidth} p{0.37\textwidth}@{}}
\toprule
Question & Mark & Awarded for & Feedback \\
\midrule
\endfirsthead
\toprule
Question & Mark & Awarded for & Feedback \\
\midrule
\endhead
1(a) & \markcell{2}{2} & Correctly identifies a northern/coastal distribution. & This gives two valid distribution statements. \\
1(b) & \markcell{4}{4} & Two valid physical reasons: poor/fertile land and lack of water access, both linked to food production or habitability. & Clear cause-effect structure. \\
1(c) & \markcell{4}{4} & Cultural reason links changing attitudes/career priorities to delayed childbearing; economic reason explains high child costs and reduced economic return from children. & Both reasons are developed enough for full marks. \\
\midrule
2(a) & \markcell{2}{2} & Defines albedo as reflected sunlight and adds the 0--1 scale / terrestrial surface idea. & Complete definition. \\
2(b)(i) & \markcell{2}{2} & Explains biomes shifting toward higher latitudes as climates warm. & Good spatial change plus causal detail. \\
2(b)(ii) & \markcell{2}{2} & Explains animal migration shifting to match changing biomes/climates. & Full credit, though a named species or specific route would make the point stronger. \\
2(c) & \markcell{4}{4} & Explains increased hurricane/cyclone hazard and vector-borne disease risk, with health consequences. & Strong, relevant examples. \\
\midrule
3(a) & \markcell{1}{2} & Identifies the rise from 2005 to 2009 and the fall from 2009 to 2014. & The markscheme requires quantification for full marks. Add approximate values from the graph. \\
3(b) & \markcell{4}{4} & Explains lack of infrastructure investment and disruption/contamination from flooding or hazards. & Both reasons are valid and developed. \\
3(c) & \markcell{4}{4} & Economic advantage: jobs/remanufacturing; environmental advantage: less extraction/mining and reduced environmental damage. & Clear distinction between economic and environmental benefits. \\
\midrule
4(a)(i) & \markcell{1}{1} & Gives 28.1 as the range difference. & Accepted by the markscheme. \\
4(a)(ii) & \markcell{1}{1} & Identifies iron. & Correct. \\
4(b) & \markcell{2}{2} & Identifies the negative relationship and notices the lowest-income anomaly. & Full credit. Adding figures such as 1.6\%, 23\%, and 15\% would make the answer more secure, but the required relationship and anomaly are present. \\
4(c) & \markcell{5}{6} & Uses evidence on hazardous landfill waste, low collection rates, landfill disposal, pollution impacts, and the value of recoverable materials. & The answer earns several supported points on both sides. The final mark was lost because there is no explicit overall appraisal weighing the infographic as a whole. \\
\bottomrule
\end{longtable}

\section*{Teaching Feedback}

\textbf{What is working well}
\begin{itemize}
  \item Explanations usually follow a clear chain: reason, process, result.
  \item Definitions and short-answer command terms are handled accurately.
  \item The student can separate social, economic, physical, and environmental factors.
  \item The climate and circular economy answers show strong conceptual understanding.
\end{itemize}

\textbf{Main improvement targets}
\begin{itemize}
  \item For graph trend questions, include figures even when the trend is obvious.
  \item In six-mark ``to what extent'' answers, end with a clear weighing-up judgement.
  \item In resource evaluation answers, keep linking each paragraph to specific evidence from the resource.
\end{itemize}

\textbf{Suggested next practice}
\begin{itemize}
  \item Rewrite Question 3(a) in one sentence with two numerical values.
  \item Annotate Question 4(b) with two percentages from the table to make the full-mark relationship more secure.
  \item Rewrite Question 4(c) using a balanced structure: two points showing e-waste is global, one point showing uneven responsibility or benefits, and one final judgement.
\end{itemize}

\section*{Model Improvements}

\textbf{Question 3(a) example structure:}
The slum population percentage increased from about \emph{[read value]} in 2005 to about \emph{[read value]} in 2009, before decreasing slightly by 2014; overall, the 2014 value remained higher than the 2005 value.

\textbf{Question 4(b) example structure:}
There is a mostly negative relationship: the highest GNI group has only 1.6\% annual EEE consumption growth, while the low GNI group has 23\%. The lowest GNI group is an anomaly because its growth rate is lower, at 15\%, than the low-income group.

\textbf{Question 4(c) example structure:}
E-waste is a global problem because the map shows international flows between sending and receiving regions, and the table shows e-waste generation across all continents. It is also dangerous because e-waste is only 2\% of solid waste but contributes 70\% of hazardous landfill waste. However, the problem is uneven: high-income countries generate much more per person, while some receiving regions may gain jobs and recover valuable materials worth EUR 55 billion. Overall, the infographic supports the view to a large extent, but it also shows that responsibility and impacts are unevenly distributed.

\section*{Notes}

\begin{itemize}
  \item Page 1 was supplied as both JPG and PNG; the files appear to be duplicate content, so it was marked once.
  \item This report records marking judgements and teaching feedback only; it does not reproduce the exam paper or markscheme.
\end{itemize}

\end{document}
